\documentclass[12pt]{article}

% === Packages ===
\usepackage[margin=1in]{geometry}
\usepackage{amsmath, amssymb}
\usepackage{setspace}
\usepackage{titlesec}
\usepackage{enumitem}
\usepackage{helvet}
\renewcommand{\familydefault}{\sfdefault} % modern sans-serif look

% === Formatting tweaks ===
\setstretch{1.15}
\setlist[enumerate]{itemsep=0.4em, topsep=0.2em}
\titleformat{\section}{\Large\bfseries\sffamily}{\thesection.}{1em}{}
\titleformat{\subsection}{\large\bfseries\sffamily}{\thesubsection}{1em}{}

% === Title ===
\title{\vspace{-2em}\textbf{Applied Mathematics Oral Exam Questions and Model Answers}\\
\large From Euler to Navier--Stokes}
\author{Prepared for Graduate Oral Examination}
\date{}

\begin{document}
\section*{Sample Oral Exam Questions: Immersed Boundary Method}

\subsection*{Conceptual Foundations}

\begin{enumerate}
\item \textbf{What is the immersed boundary (IB) method, and what class of problems was it originally designed to study?}

\textit{Sample answer:}
The immersed boundary method is a computational framework for simulating fluid--structure interaction (FSI) problems. It was originally developed by Charles Peskin in the 1970s to study blood flow in the heart, particularly the interaction between elastic heart valves and viscous blood. The method couples a Lagrangian representation of the structure with an Eulerian representation of the fluid.

\item \textbf{Describe the Eulerian--Lagrangian coupling used in the immersed boundary method.}

\textit{Sample answer:}
In the IB method, the fluid variables (velocity and pressure) are defined on a fixed Eulerian grid, while the structure is represented by Lagrangian points or elements. Forces generated by the structure are computed in Lagrangian coordinates and then spread to the Eulerian grid using a regularized delta function. The resulting fluid velocity is interpolated back to the Lagrangian points to update the motion of the structure.

\item \textbf{What does it mean that boundary conditions are enforced “weakly” in the immersed boundary method?}

\textit{Sample answer:}
Weak enforcement means that boundary conditions, such as the no-slip condition, are not imposed pointwise at the fluid--structure interface. Instead, they are satisfied in an averaged or integral sense through force spreading and velocity interpolation. As a result, velocity is continuous across the interface, while pressure and stress may be discontinuous.
\end{enumerate}

\subsection*{Mathematical and Numerical Structure}

\begin{enumerate}[resume]
\item \textbf{Write down the generic form of the force coupling in the immersed boundary method.}

\textit{Sample answer:}
If $\mathbf{F}(\mathbf{X},t)$ denotes the Lagrangian force density on the structure and $\mathbf{f}(\mathbf{x},t)$ the Eulerian force density, then
\[
\mathbf{f}(\mathbf{x},t) = \int \mathbf{F}(\mathbf{X},t)\,\delta_h(\mathbf{x}-\mathbf{X})\,d\mathbf{X},
\]
where $\delta_h$ is a regularized delta function. The structure moves according to
\[
\frac{d\mathbf{X}}{dt} = \int \mathbf{u}(\mathbf{x},t)\,\delta_h(\mathbf{x}-\mathbf{X})\,d\mathbf{x}.
\]

\item \textbf{What factors control the accuracy of an immersed boundary simulation?}

\textit{Sample answer:}
Accuracy depends on the Eulerian grid spacing, the Lagrangian mesh spacing, and the properties of the regularized delta function, including its smoothness and support. Consistency between Eulerian and Lagrangian resolutions is critical. As both grids are refined, the regularized formulation converges to the continuum IB equations.

\item \textbf{Why can stiff elastic forces lead to severe time-step restrictions in explicit IB schemes?}

\textit{Sample answer:}
Stiff elastic forces produce fast structural dynamics that require very small time steps for numerical stability when treated explicitly. This can make simulations computationally expensive. Many modern IB extensions address this issue using implicit or semi-implicit coupling strategies.
\end{enumerate}

\subsection*{Advantages, Limitations, and Failure Modes}

\begin{enumerate}[resume]
\item \textbf{List two major advantages and two major disadvantages of the immersed boundary method.}

\textit{Sample answer:}
Advantages include the ability to handle large structural deformations and complex, moving geometries without remeshing, and a unified framework for modeling fibers, surfaces, and solids. Disadvantages include numerical smearing of the interface, potential slip at the boundary, and restrictive time-step constraints for stiff structures.

\item \textbf{Give an example of a common failure mode in IB simulations and explain its cause.}

\textit{Sample answer:}
One common failure mode is force leakage, which occurs when the Lagrangian mesh is too coarse relative to the Eulerian grid. This leads to poor force transmission between the structure and fluid. Another failure mode is excessive stiffness when the Lagrangian spacing is too fine, resulting in severe time-step restrictions.
\end{enumerate}

\subsection*{Comparison with Other FSI Methods}

\begin{enumerate}[resume]
\item \textbf{Compare the immersed boundary method with grid-fitted (ALE or body-fitted) FSI methods.}

\textit{Sample answer:}
Grid-fitted methods enforce boundary conditions sharply and provide high accuracy near interfaces, but they require moving meshes and frequent remeshing for large deformations. The immersed boundary method uses a fixed Cartesian grid, avoids remeshing, and is more robust for large, complex motions, at the cost of reduced interface sharpness.

\item \textbf{How does the immersed interface method (IIM) differ from the immersed boundary method?}

\textit{Sample answer:}
The immersed interface method enforces exact jump conditions across interfaces by modifying finite difference stencils near the boundary, resulting in a sharp interface representation. In contrast, the immersed boundary method spreads forces using regularized delta functions, producing a diffuse interface but greater robustness for large deformations and complex structural behavior.
\end{enumerate}

\subsection*{Applications and Practical Considerations}

\begin{enumerate}[resume]
\item \textbf{Give two biological and two engineering applications of the immersed boundary method.}

\textit{Sample answer:}
Biological applications include blood flow in the heart and swimming of organisms such as fish and jellyfish. Engineering applications include flapping flags, flow-induced vibrations of flexible beams, and soft robotic systems.

\item \textbf{Why is the availability of open-source IB software an important practical advantage?}

\textit{Sample answer:}
Open-source IB frameworks, such as IBAMR and IB2d, lower the barrier to entry for new researchers, improve reproducibility and transparency, and allow users to focus on scientific questions rather than implementing complex solvers from scratch. This is especially important for large-scale 3D simulations.
\end{enumerate}

\end{document}